\subsection{Mode competition and finite-amplitude gain in PING}\label{supp:ping-coupling-results}
In NMM2 PING, $J$ is the mean inter-neuron QIF synaptic weight.
Sweeping $J$ at fixed background drive toward the gamma Hopf reveals two
competing modes (Fig.~\ref{fig:nmm2_jcurve}). A second mode near $55$~Hz carries
the linear-response peak over $14.6\lesssim J\lesssim16.2$; the critical mode
carries it at lower couplings and again near the Hopf, where its frequency is
about $50$~Hz. In that near-Hopf interval,
$0.0014\lesssim\gamma\lesssim0.0081$~ms$^{-1}$, the linear gain scales as
$1/\gamma$ with a fitted exponent of $1.02$ over a $5.8$-fold damping range.
The full-sweep exponent is $0.71$, so a single-mode law does not describe the
whole curve. Coupling-dependent sensitivity \cite{clusella2023} therefore
includes changes in which mode carries the peak.

At $A_f=8$, the simulated response exceeds the normalized linear reference by
up to $48\%$ in the amplitude check, then saturates. At a quarter of that
amplitude ($A_f=2$, one sixteenth of the power), the departures stay within
$11\%$ at the nine tested couplings. The driven endpoint remains attracting under both waveforms
(\S\ref{app:ping-stability}).
The alpha JR and gamma NMM2 circuits share the $1/\gamma$ gain scaling of their critical mode despite their different
mixers, with amplification adjustable through coupling.
